% AI-written, AI cross-reviewed; no human sign-off; not a claim that RH is proved
% Verbatim sections of local blueprint v13. See README.md for macros, edition and source spans.

% blueprint/ch_P.tex:178-649

% Start of atlas edition annotation. The following block is outside the original excerpt.
\paragraph{Status of named inputs in this reading.}
External theorem numbers in the original excerpt are not verified here unless a separate edition-specific source record documents the comparison.
% End of atlas edition annotation.

\subsection{Which finite characters have to be counted?}

Fix this $M$ for the next three steps and write $r=r_2(M)$,
$\kappa=1+\ell^c$, $a=\kappa^E$. At level $\nu\ge c$ write $q=\ell^\nu$.
The relevant finite Galois layer $F/M$ has a primitive $q$-th root $\zeta$
and a generator $\gamma$ acting by $\zeta\mapsto\zeta^\kappa$ of order
$h=\ell^{\nu-c}$. Define $\mathcal H_q$ to consist of functions on the
nonzero prime ideals $P$ of $\mathcal O_F$ with values in $\mathbf Z/q$
satisfying
\[
 \sum_P \operatorname{ord}_P(\alpha)\phi(P)=0\quad(\alpha\in F^\times),
 \qquad \phi(\gamma^{-1}P)=a\phi(P).
\]
The first condition factors the function on divisors through the ideal
class group. The inverse in the second records the action used for
characters. Write
\[
 N_a(\alpha)=\prod_{j=0}^{h-1} \gamma^j(\alpha)^{a^j}.
\]
Call $\alpha$ eligible if for a finitely supported integral divisor $u$
one has $\operatorname{ord}_P(\alpha)=u(\gamma P)-a u(P)$ for every $P$.
All these conventions, including the placement of the powers, are
\nolinkurl{Lay}, \nolinkurl{DivDel}, \nolinkurl{nrm} and
\nolinkurl{EqHom} (\path{ClassCounting/HasVal.lean:120--139}).

\begin{proposition}[Separation by evaluations]\label{prop:P-k-few}
\lean{LateralTheory.Q.SFewOfCut.sFewOf}
\uses{prop:P-k-layer}
Set
\[
 b=2E+(2E+1)(c+2E+1),\quad N_0=4E+6,\qquad
 T_i=\left|H_i(\mathrm{GL}_{N_0}(\mathcal O_M),\mathbf Z)_{\rm tors}\right|.
\]
If a finite set $S$ of eligible $\alpha$ has pairwise distinct
$N_a(\alpha)^{\ell^b}$ in $F^\times/(F^\times)^q$, then
$|S|\le q^r T_{2E+1}T_{2E}$. The torsion here is the entire integral
torsion group (\nolinkurl{torsHGL}, \path{GeneratorCountTuples/TorsBoundPow.lean:13}).
\end{proposition}
\begin{proof}
We specify the chain and detection inputs at the point they enter.

\emph{K2, cycle construction and evaluation.} Section~\ref{sec:P-r3-K247} constructs the relation and interpolating multiplicities. The scalar evaluator, integral transport and uniform reduction are proved in the K2 argument at \ref{sec:P-r4-K2}.
Put $m=2E+1$. For each eligible $\alpha$ with divisor $u$ there
exist a degree-$m$ cycle $y_\alpha$ modulo $q$ in
$\mathrm{GL}_{N_\alpha}(\mathcal O_M)$ and a nonzero
$b_\alpha\in\mathcal O_F$. For every reduction $\mathcal O_F\to k$ to a
finite field that is surjective on $\mathcal O_M$ and preserves
$b_\alpha\ne0$, its evaluation by any admissible additive functional
with index character $v:k^\times\to\mathbf Z/q$ is
\[
  P_E\,v\big(\overline{N_a(\alpha)}\big),\qquad
  P_E=\prod_{\substack{0\le i<2E+1\\ i\ne E}}
                   (\kappa^E-\kappa^i).
\]
Here the bar means a residue represented by an integral numerator and
denominator with nonzero denominator. At a prime away from $\ell$ the
root $\zeta$ stays primitive. For a chosen $v$ one can choose such a
functional common to all the $y_\alpha$; its pullback is a family
vanishing on boundaries, stable under enlargement of the matrix size and
additive under block repetition on cycles. These are exactly the uses of
\nolinkurl{SClass}, \nolinkurl{STheta} and \nolinkurl{SPhi}
(\path{LemmaKOfSteps.lean:54--69,84--96}), with their proofs instantiated in
\nolinkurl{lemmaK} (\path{ClassCounting/LemmaKOfLitN.lean:28--31}).
The degree is $2E+1$ because the input there is indexed by $e=E-1$.
In particular it includes $E=1$.

Exclude primes dividing $\ell\prod_{\alpha\in S}b_\alpha$; the product
is nonzero even for the empty set. A primitive $\ell^\nu$-th root remains
primitive at each remaining prime: if its $\ell^{\nu-1}$-st power
reduced to $1$, the geometric sum for a primitive $\ell$-th root would
give $\ell=0$ in the residue field. Its multiplicative group is cyclic of
order divisible by $q$. Taking an exponent modulo $q$ in a cyclic
generator gives an index character $v$ with kernel the $q$-th powers
(\nolinkurl{prim_map}, \nolinkurl{exists_index},
\path{SFewOfCut.lean:85--117}).

Use all the allowed primes $P$ at which $\mathcal O_M\to k(P)$ is
surjective. Assign to $\alpha$ the vector of evaluations of
$\ell^{2E}y_\alpha$ in $\prod_P\mathbf Z/q$. If two vectors agree, then
$\ell^{2E}P_E$ kills the difference of the index characters of the two
twisted norms at every such $P$. In $\mathbf Z/q$ the multiple
$\ell^{(2E+1)(c+2E+1)}$ belongs to the ideal generated by $P_E$.
For this estimate each factor is a unit times
$\kappa^{|E-i|}-1$, with $1\le |E-i|<2E+1<\ell^c$;
the $\ell$-adic valuation is $c+v_\ell(|E-i|)$ (at $2$, $c\ge2$).
There are at most $2E+1$ factors. Removing the prime-to-$\ell$ part
modulo $q$ gives the asserted ideal containment even when the bound
already exceeds $\nu$ (\nolinkurl{pe_dvd}, \path{SFewOfCut.lean:67--83}).
The ratio of the two $\ell^b$-powered norms consequently reduces to a
$q$-th power at every allowed prime. To pass from residues to $F$, choose a $q$-th root $x$ of this ratio
$\beta$ in an algebraic closure. For $g\in G_F$ the set of $\sigma$ with
$\sigma x=gx$ is a nonempty open coset. The density input \emph{K3, proved in Proposition~\ref{prop:P-k-cheb}} says that the Frobenius elements with degree-one prime in $F$
over $\mathbf Q$, allowing the choice of a prolongation, meet every such
open subset even after excluding finitely many primes
(\nolinkurl{ChebotarevDegreeOne.dense}, \path{SDetectCut.lean:111--126}).
The residue field at the chosen prime has prime order. Its prime subfield
already fills it, so $\mathcal O_M$ surjects onto it. By the test, $\beta$
reduces to some unit $\bar t^q$. Write $t$ for an integral representative.
The factorization of $x^q-t^q$ over $\mu_q$ gives
$x\equiv\zeta^i t$ at the chosen prolongation for some $i$. Indeed,
the product of the linear factors has positive additive valuation, so
at least one factor does. Thus $x$ is a unit with residue in the base
residue field. Its Frobenius conjugate $\sigma x$ has the same residue,
and $\sigma x/x$ is a $q$-th root of unity reducing to $1$. Away from
$q$ that root must equal $1$: for a nontrivial root the geometric sum
of length $q$ is zero and its reduction would give $q=0$.
Consequently $gx=\sigma x=x$ for every $g$. The fixed-field theorem gives
$x\in F^\times$ and $\beta=x^q$. This proves the implication
\nolinkurl{SDetect} (\path{LemmaKOfSteps.lean:109--115}) from the
degree-one density assertion (\nolinkurl{frob_fix}, \nolinkurl{goal},
\path{SDetectCut.lean:43--99,101--142}). Distinctness now forces the two
members of $S$ to coincide.

Finally the vectors form evaluations of a \emph{single} admissible
family with values in the product module. Its operators, and the
identities for boundaries, stabilization and repeated blocks, are
defined coordinate by coordinate. The required bound on these images, K4, is as follows: for even $k$
with $k+1<r$, any such family with values in a $\mathbf Z/q$-module and
any finite set of values $\ell^k\Phi(y)$ on degree-$(k+1)$ cycles in
arbitrary matrix sizes, its cardinality is at most
$q^r$ times the torsion orders in degrees $k+1,k$ in rank
$2(k+1)+4$. See \nolinkurl{SUpper}
(\path{LemmaKOfSteps.lean:100--107}) and \nolinkurl{sUpperOf}
(\path{SUpperOf.lean:21--66}), whose reduction uses the Bockstein and
the rank bound \nolinkurl{BOdd} (\path{LemmaKOfSteps.lean:11--15}).
All matrix sizes occurring in a finite set can first be stabilized into
one common rank. With $k=2E$ the stable rank is $N_0$. Injectivity just
proved gives the stated inequality (\nolinkurl{sFewOf},
\path{SFewOfCut.lean:153--211}).
Proposition~\ref{prop:P-k-bock} derives K4 from a stable torsion comparison and a rank bound for the indicated classes. \emph{Status: complete modulo K4.}
\end{proof}

\begin{proposition}[The degree-one Frobenius test]\label{prop:P-k-cheb}
\lean{GlobalTheory.Published.ChebotarevDegreeOne.dense}
\uses{def:S-lemmaK}
For a number field $F$, an algebraic closure $\Omega/F$ and a finite
set $S$ of finite places, the arithmetic Frobenius lifts at absolute
degree-one primes outside $S$, allowing every choice of prolongation,
form a dense subset of $G_F$. This is the specialization of
\nolinkurl{ChebotarevDegreeOne.dense} to \nolinkurl{chebotarevDegreeOne'}
(\path{ChebotarevDegreeOneIffCheb.lean:12--18,51--62}).
\end{proposition}
\begin{proof}
\par\noindent\textbf{Alternative argument via finite Chebotarev.} 
We give the specialization to the coset used there. Fix $x$ and
$g\in G_F$. In the chosen algebraic closure of $F$, take a finite
extension $N$ Galois over $\mathbf Q$ containing $F$ and $x$.
For instance take the normal closure over $\mathbf Q$ of $F(x)$.
Then $\tau=g|_N$ belongs to $\operatorname{Gal}(N/\mathbf Q)$ and fixes
$F$ pointwise. The Chebotarev theorem applies because $N/\mathbf Q$
is finite and Galois. The conjugacy class of $\tau$ has positive density
among unramified rational primes. We may discard the finite set of
rational primes below the excluded prime ideals of $\mathcal O_F$.
See Milne \cite[Theorem~8.31]{P-Milne-Cheb}.

Select a remaining rational prime $p$ in the specified class. There
is a prime $\mathfrak P$ of $N$ over $p$ whose actual arithmetic
Frobenius is $\tau$: conjugating a prime conjugates its Frobenius,
so start with any prime over $p$ and translate it by the desired
conjugating automorphism. Put $P=\mathfrak P\cap\mathcal O_F$. For
each $z\in\mathcal O_F$ the identities
\[
   z=\tau z\equiv z^p\pmod{\mathfrak P}
\]
show that every residue in $k(P)$ is fixed by $p$-power Frobenius.
Hence $k(P)=\mathbf F_p$. The same $\tau$ is the arithmetic Frobenius
of $\mathfrak P/P$ for $N/F$, since its action on the residue field
of $N$ is raising to $|k(P)|=p$. In particular its action on $x$ is
$gx$. This proves exactly the specialization needed to test the
$q$-th power, including the surjectivity of the reduction of
$\mathcal O_M$: a unital map to $\mathbf F_p$ already hits that
field through its rational integers.

For the open-set formulation, a neighborhood of $g$ in $G_F$
contains the coset determined by its action on a finite Galois
extension of $F$. Include that entire extension in $N$ in place of
$F(x)$. Prolong $\mathfrak P$ to the algebraic closure and choose an
arithmetic Frobenius lift in the decomposition group over $F$.
It restricts to $\tau$ on $N$ because the prime is unramified there,
so it lies in the neighborhood.

\par\noindent\textbf{Lean route via a cyclotomic extension.} 
The formal proof making the density step calls \nolinkurl{ChebotarevDegreeOne.dense} on the nonempty
coset (\path{SDetectCut.lean:111--128}); its source
\nolinkurl{chebotarevDegreeOne'} (\path{InertiaEigenvalues/ChebotarevDegreeOne.lean:17--64})
lifts the prescribed automorphism of the finite Galois extension,
selects an auxiliary prime and a crossing field whose action on roots
of unity detects that automorphism. It excludes ramified primes, finds a
prime of prescribed prime absolute norm in the crossing field and restricts
its Frobenius to $F$. The final roots-of-unity calculation recovers the
given element of the finite quotient. The finite Chebotarev argument
above is an independent hand proof of K3. The internal construction
supplies K3 to \nolinkurl{SDetectCut.goal}
(\path{SDetectCut.lean:101--128}).
\emph{Status: K3 complete by the independent published route.}
\end{proof}
%BIB P-Milne-Cheb | J. S. Milne | Algebraic Number Theory | University of Michigan course notes, version 2.10 (1998) | https://www.jmilne.org/math/CourseNotes/ANT210.pdf | Theorem 8.31, positive density for every Frobenius conjugacy class in a finite Galois extension of number fields.

\subsection{Why the two consecutive homology torsion groups occur}
\label{sec:P-k-bock}

\begin{proposition}[The integral-cycle obstruction and the rank estimate]
\label{prop:P-k-bock}
\lean{UpperTheory.V.ZK.selmer_upper_bound_of_inputs}
\uses{def:S-lemmaK}
In the evaluation bound of Proposition~\ref{prop:P-k-few}, the remaining
homological input K4 can be reduced to the following two assertions.
Here $M$ is totally complex, $k$ is even, $k+1<r_2(M)$, and
$N_0=2(k+1)+4$.
\begin{enumerate}
\item In degree $i=k,k+1$, for $H_i(\mathrm{GL}_d(\mathcal O_M),\mathbf Z)$
the torsion is finite and has the same order at $d=N_0$ and all $d\ge N_0$.
This is the specialization of \nolinkurl{TorsStab},
\path{TorsionRankCounts/RankCount.lean:28--31}. In the project it follows
from integral homology stabilization at $d\ge 2i+4$ and finiteness in the
initial rank (\nolinkurl{torsStab}, \path{TorsStab.lean:16--39}).
\item Start at any $d\ge N_0$. For $x\in H_{k+1}(\mathrm{GL}_d(\mathcal O_M),\mathbf Z)$
define $x_0=x$, and obtain $x_{j+1}$ by taking the homology image of $x_j$
under $\ell$ identical diagonal blocks, minus $\ell^{j+2}$ times its
ordinary stabilized homology image in that same rank. The span over
$\mathbf Z$ of all $x_k$, at rank $\ell^k d$, has rank at most $r_2(M)$.
This is the needed case of \nolinkurl{BOdd},
\path{LemmaKOfSteps.lean:11--15}, with the two maps in
\nolinkurl{psi}, \nolinkurl{zSeq}, \path{ClassCounting/HasVal.lean:45--60}.
\end{enumerate}
\end{proposition}
\begin{proof}
Let $S$ be a finite set of evaluations of the allowed shape.
Each $s\in S$ has a cycle representing it in some matrix rank $d_s$.
Take $d=\max(N_0,\max_{s\in S}d_s)$, with the maximum of the empty
collection taken as zero. Inserting identity blocks induces a chain map,
and the evaluation on cycles is stable under this operation. We can thus
use a single group $G=\mathrm{GL}_d(\mathcal O_M)$ for the entire count.
Put $q=\ell^\nu$ and let $B$ be the common target of the evaluations, so
$qB=0$. Write $g(y)=\ell^k\Phi_d(y)$ for a degree-$(k+1)$ cycle
modulo $q$. We examine the cost of lifting $y$ to an integral cycle
before estimating the possible values of $g$. The formal
\nolinkurl{selmer_upper_bound_of_inputs} takes six intermediate results
and yields $\texttt{SChain}\Rightarrow\texttt{BOdd}\Rightarrow\texttt{SUpper}$
(\path{SUpperOf.lean:24--45}). In the explicit bar model the cycle-to-homology
map is onto, with kernel the boundaries and functoriality under group
maps (\nolinkurl{Up.bridge}, \path{TorsionRankCounts/Bridge.lean:5--56}).
On tuple generators the integral bar boundary has the usual alternating
faces, so coefficientwise reduction commutes with the differential and
with group maps (\nolinkurl{Up.rqFacts}, \path{TorsionRankCounts/RqFacts.lean:5--38}).

\smallskip\noindent\emph{The obstruction in degree $k$.}
Let $C_\bullet$ be the ordinary integral bar complex. Its terms are
free abelian on tuples of group elements, so reduction modulo $q$
is onto and has kernel exactly $qC_i$ in degree $i$. Lift $y$ to
$\widetilde y\in C_{k+1}$. Since $y$ is a cycle,
$\partial\widetilde y=q w$. Then $q\partial w=0$ by $\partial^2=0$,
and torsion-freeness gives $\partial w=0$. Changing the lift by $q z$
changes $w$ by $\partial z$. We have therefore constructed an additive map
\[
 \beta:Z_{k+1}(C_\bullet/qC_\bullet)\longrightarrow H_k(C_\bullet)[q],
              \qquad y\longmapsto[w] .
\]
Its values are killed by $q$ since $q w$ is a boundary. If $[w]=0$,
choose $t$ with $w=\partial t$ and replace $\widetilde y$ by
$\widetilde y-q t$. This is an integral cycle with reduction exactly
$y$, including equality of chains. Conversely, an integral cycle lift
gives $[w]=0$. The construction is additive because we may add the lifts.
These are \nolinkurl{bock_abs} and \nolinkurl{bock}
(\path{TorsionRankCounts/BockAbs.lean:8--70,92--105}).
For $k=0$, the boundary $C_1\to C_0$ is zero for trivial coefficients,
and each degree-one chain already has an integral cycle lift
(\nolinkurl{bd_zero}, \path{BockAbs.lean:82--90}).

For each $s\in S$ choose one $y_s$ with $g(y_s)=s$. Partition this
finite set by $\beta(y_s)$. In a nonempty part choose $s_0$. The map
$s\mapsto s-s_0$ is injective on the part and
$s-s_0=g(y_s-y_{s_0})$, where $\beta(y_s-y_{s_0})=0$. A uniform bound
for finite sets of evaluations of liftable cycles therefore bounds each
part. The number of nonempty parts is at most the order of the entire
torsion subgroup of $H_k(G,\mathbf Z)$. This accounts for the degree-$k$
factor even when the Bockstein is far from surjective
(\nolinkurl{fibCount}, \path{TorsionRankCounts/FibCount.lean:5--27}).

\smallskip\noindent\emph{Evaluating the integral lifts.}
If $y$ has an integral cycle lift $z$, evaluating its reduction depends
on $[z]\in H_{k+1}(G,\mathbf Z)$: reducing an integral boundary gives
a boundary, and $\Phi_d$ vanishes on boundaries. Denote the resulting
homomorphism on integral homology by $\phi_d$. The equalities for
stabilization and block repetition pass to these maps by choosing a cycle
representative. They say, respectively, that ordinary stabilization leaves
the value unchanged and that $j$ identical diagonal blocks multiply the
value by $j$ (\nolinkurl{descend}, \path{TorsionRankCounts/Descend.lean:8--55}).

Let $d_j=\ell^j d$ and take $x=[z]$ in the recurrence in the statement.
Its two terms are in the same integral homology group, with evaluations
$\ell\phi_{d_j}(x_j)$ and $\ell^{j+2}\phi_{d_j}(x_j)$.
Inductively one obtains
\[
 \phi_{d_t}(x_t)=\ell^t A_t\,\phi_d(x),\qquad
 A_t=\prod_{0\le j<t}(1-\ell^{j+1})\equiv1\pmod\ell .
\]
In particular $A_0=1$ and all $A_t$ are invertible modulo $q$. If
$A_t=1-\ell c_t$ then $\sum_{i=0}^{\nu-1}(\ell c_t)^i$ is an inverse
modulo $\ell^\nu$. Multiplication by it in $B$ is legitimate because
$qB=0$. Taking $t=k$ expresses our value $g(y)=\ell^k\phi_d(x)$ as
$\phi_{d_k}$ of an integral multiple of $x_k$. No division by $\ell$
has occurred. This distinction matters for an evaluation target with
elements of order $\ell$. The recurrence and its inversion are
\nolinkurl{phi_zSeq}, \nolinkurl{zVal}
(\path{TorsionRankCounts/PhiZSeq.lean:5--35}).

\smallskip\noindent\emph{Counting with the rank hypothesis.}
Consider any finite set of values from the liftable cycles.
For each value take the integral multiple just constructed and generate
a subgroup $J'$ of $H_{k+1}(\mathrm{GL}_{d_k}(\mathcal O_M),\mathbf Z)$
by these finitely many preimages. By the second hypothesis it has rank
at most $r_2(M)$. The quotient of $J'$ by its torsion is a finitely
generated free abelian group. Pick lifts $v_1,\ldots,v_t$ of a basis,
$t\le r_2(M)$. An element of $J'$ is a sum $\sum a_i v_i+w$ with
$w$ torsion. In $B$ we may replace each $a_i$ independently by its
residue in $\{0,\ldots,q-1\}$. Thus the possible values in this finite
set number at most
\[
  q^{r_2(M)}
     \left|H_{k+1}(\mathrm{GL}_{d_k}(\mathcal O_M),\mathbf Z)_{\rm tors}\right|.
\]
Only $J'$ had to be finitely generated. A rational rank bound on an
arbitrary ambient subgroup alone would not provide integral generators.
This is \nolinkurl{rankCount_gen}, \nolinkurl{rankCount}
(\path{TorsionRankCounts/RankCountAlt.lean:5--23,35--78}).

Use the first hypothesis for the torsion factor in degree $k+1$ at $d_k$
and for the Bockstein factor in degree $k$ at $d$. Both ranks are at
least $N_0$. Multiplying the last bound by the number of parts gives
exactly $q^{r_2(M)}T_{k+1}T_k$. In Proposition~\ref{prop:P-k-few} the
two assertions of K4, proved in Section~\ref{sec:P-r5-k4}, are required at $k=2E$. In particular the
rank bound concerns the span of the recursively constructed classes in
degree $2E+1$; a bound for the whole homology group has not been used.
\emph{Status: complete using the invariant-form comparisons of Section~\ref{sec:P-r5-k4}; Proposition~\ref{prop:P-r4-k7} supplies the torsion finiteness.}
\end{proof}

\subsection{Removing the growing power of the level}

\begin{proposition}[Divisors, units and finite subsets]\label{prop:P-k-count}
\lean{UpperTheory.V.ZK.sCount}
\uses{prop:P-k-few}
With $v=v_\ell(E)$, $g=2n'+\lfloor\log_2 h_M\rfloor+2$, where $h_M$ is
the class number, and $X=(b+v)g+c+v$, one has
\[
 |\mathcal H_q|\le B_*:=T_{2E+1}T_{2E}\ell^X.
\]
Moreover $\Sigma(M)=s$ for an integer $s\ge0$ and $\ell^s\le B_*$.
The first assertion combines \nolinkurl{sCount} with separation
and cancellation (\nolinkurl{sNum}, \path{NumThree.lean:221--241}).
The second is \nolinkurl{PositivityTheory.QuaternionShaBounds.sSha}
(\path{OrderOfOneAddPrimePowZmodPrimePow.lean:490};
\nolinkurl{SSha}, \path{ClassCounting/LemmaKOfSteps.lean:143--147}).
\end{proposition}
\begin{proof}
In additive notation take $Y=F^\times$, $I=\operatorname{Div}(\mathcal O_F)$,
with principal divisor map $\partial:Y\to I$. On these groups write
$\delta$ for the action of $\gamma^{-1}$ minus $a$ on $Y$.
On divisors the same endomorphism sends $u(P)$ to $u(\gamma P)-a u(P)$,
so it commutes with $\partial$ (\nolinkurl{IsDat}, \path{TwoLe.lean:56--61}).
Let $U=\partial^{-1}(\delta I)$. Form
\[
 W=U/(\delta Y+Y_{\rm tors}),\qquad
 Q_{\rm cl}=I/(\partial Y+\delta I),\qquad
 Q_{\rm u}=\ker\partial/(\delta\ker\partial+Y_{\rm tors}).
\]
These are \nolinkurl{Dat.W}, \nolinkurl{Dat.ClQ}, \nolinkurl{Dat.EQ}
(\path{SelmerOrderCounts/TwoLe.lean:34--47}).
Lemma~\ref{lem:P-r3-k5} proves the following arithmetic properties.
All three groups are finite,
$|W|=|Q_{\rm u}|\,|Q_{\rm cl}|$,
$|\mathcal H_q|\le\ell^{v_\ell(|Q_{\rm cl}|)}$ with $\mathcal H_q$ finite,
and $|W[\ell]|\le\ell^g$. The powered norm gives a homomorphism $f$ on
$W$ whose image can be represented by a finite family as in
Proposition~\ref{prop:P-k-few}. Its kernel is killed by
$m\ell^{b+v}$ with $(m,\ell)=1$. See \nolinkurl{sCount}
(\path{SelmerOrderCounts.lean:20--71}) for \nolinkurl{cardW},
\nolinkurl{eqHomLe}, \nolinkurl{kerF}, \nolinkurl{fewLe} and
\nolinkurl{selInj} in this application. In particular the condition
on $W[\ell]$ is a condition at the base $M$, although the group $W$ was
constructed using $F$.

The unit factor supplies $q^r$. The project takes the logarithmic
unit lattice modulo roots of unity and appends one trivial integral
coordinate to fill the constant direction. On the full real space of
archimedean places, $\gamma$ permutes the places in $r$ cycles of
length $h$ because $M$ is totally complex. On one cycle the absolute determinant
of this permutation minus $a$ times the identity is $a^h-1$
(the characteristic polynomial is $t^h-1$). On the extra constant line
it is $a-1$. Computing the index of the resulting lattice endomorphism
by its determinant gives
\[
 |Q_{\rm u}|(a-1)=(a^h-1)^r .
\]
Lemma~\ref{lem:P-r3-k5} identifies this logarithmic lattice with the
units modulo roots of unity. The determinant and division account for the two different
integer factors in this formula (\nolinkurl{cardEQ},
\path{SelmerOrderCounts/CardEQ.lean:98--128}). Since $\kappa=1+\ell^c$
and $h=\ell^{\nu-c}$,
\[
 v_\ell(a-1)=c+v,\qquad v_\ell(a^h-1)=\nu+v.
\]
Indeed raising a principal unit to an exponent prime to $\ell$ preserves
the valuation of its difference from $1$, and raising it to $\ell$
increases that valuation by one when $c\ge1$ for odd $\ell$ and $c\ge2$
for $\ell=2$. Both facts follow by expanding $(1+t)^j$ and comparing
the first term with subsequent terms. In particular $a^h-1=q\ell^v m$
for some $m$ prime to $\ell$, also when $\nu=c$ and $h=1$. Thus
\[
 v_\ell(|Q_{\rm u}|)+c+v=r(\nu+v)\ge r\nu .
\]
This is the equality used in \nolinkurl{sCount}
(\path{SelmerOrderCounts.lean:30--50,81--98}).

Let $K_f$ be the kernel of $f$. On its $\ell$-primary subgroup $m$ acts
invertibly, so that subgroup is killed by $\ell^{b+v}$. Each multiplication
by $\ell$ from $W[\ell^i]$ to $W[\ell^{i-1}]$ for $i\ge1$ has kernel of size at
most $|W[\ell]|$. Induction gives $|W[\ell^{b+v}]|\le\ell^{(b+v)g}$.
The index of $K_f$ in $W$ is $|\operatorname{im} f|$. Taking $\ell$-parts
yields
\[
 \ell^{v_\ell(|W|)}\le|\operatorname{im}f|\,\ell^{(b+v)g}.
\]
This is \nolinkurl{pCount} (\path{SelmerOrderCounts/PCount.lean:38--92}).
Multiplying the class-character bound by $q^r$ and inserting the unit
identity gives
\[
\begin{aligned}
 |\mathcal H_q|q^r
 &\le\ell^{v_\ell(|Q_{\rm cl}|)+v_\ell(|Q_{\rm u}|)+c+v}\\
 &\le |\operatorname{im}f|\,\ell^{(b+v)g+c+v}
 \le q^r T_{2E+1}T_{2E}\ell^X .
\end{aligned}
\]
Since $q^r>0$ it can be cancelled. In particular $X$ uses $c$, which was
fixed when $M$ was chosen, and $v_\ell(E)$. It contains no $\nu$.
This is the passage from \nolinkurl{sCount}
(\path{SelmerOrderCounts.lean:91--111}) to \nolinkurl{sNum}
(\path{NumThree.lean:221--241}).

Lemma~\ref{lem:P-r3-k6} supplies the passage back to cohomology. Restriction to the infinite
cyclotomic layer injects the degree-one localization kernel into the
equivariant continuous characters there. Each character in that image
lifts, at every sufficiently large $\nu$, to an equivariant finite-level
ideal class character killed by $\ell^\nu$. For the cyclic layer with
generator chosen above there is an injection from these latter characters
into $\mathcal H_q$. The finite-level pullback determines the original
infinite-level character. The identifications use
\nolinkurl{good_of}, \nolinkurl{ClInj}, \nolinkurl{LayEx},
\nolinkurl{ShaAsm_holds}
(\path{OrderOfOneAddPrimePowZmodPrimePow.lean:26--38,390--410,454--463}).
In the source the threshold is obtained by first extending an equivariant
character to a finite level, expressing it by pullback from a finite class
group there, and bounding the exponent of its finite image. Restriction
to deeper levels preserves the bound on its image. This explains why all
sufficiently large levels work.

Take a finite subset $T$ in the image of global restriction. A maximum
of its thresholds, $c$ and $2$ gives one level for all its members.
Choose one lift per member. If the chosen finite-level images in
$\mathcal H_q$ coincide, injectivity before restriction gives equality of
the lifts, and restriction gives equality of their members in $T$.
Thus $|T|\le B_*$. An infinite image would have a finite subset of size
$B_*+1$. The image is consequently finite, and applying the same bound
to the entire image gives its order at most $B_*$. A finite
$\mathbf Z_\ell$-module has order $\ell^s$ for its length $s$ (take
successive simple quotients). Restriction is injective, hence this $s$
is $\Sigma(M)$. This finite-subset argument is \nolinkurl{card_le} and
\nolinkurl{ShaAsm_holds} (\path{OrderOfOneAddPrimePowZmodPrimePow.lean:412--448,464--490}).
\emph{Status: complete modulo Proposition~\ref{prop:P-k-few}.}
\end{proof}

